\chapter{Project Management Plan}
\section{Team Work}
\subsection{Team Structure and Roles}
Table~\ref{tab:team} records the four members supplied in the course document. Individual role assignments, time sheets, and a signed responsibility allocation are not present in the inspected project artifacts. The report therefore describes responsibilities by work package rather than attributing undocumented work to particular students.

\begin{table}[H]\centering\small
\caption{Functional responsibilities and reviewable outputs.}\label{tab:responsibilities}
\begin{tabularx}{\textwidth}{P{3cm}YY}\toprule
\textbf{Responsibility} & \textbf{Engineering output} & \textbf{Review evidence}\\\midrule
Data preparation & Pinned source selection, normalized metadata, hierarchical chunks & Corpus manifest, source revision, chunk identities\\
Retrieval and context & Lexical/dense/graph candidates, fusion, reranking, packing & Ranking traces, retained source text, token budgets\\
Orchestration and reliability & Sufficiency decisions, bounded feedback, generation, verification & Status, round history, structured failures, regression tests\\
Evaluation and reporting & Question validation, scoring, statistics and visual analysis & Frozen run manifest, score artifacts, CSV/JSON exports\\\bottomrule
\end{tabularx}
\end{table}
These responsibilities provide a practical basis for peer review. Changes that cross boundaries should be reviewed against the relevant contracts: for example, a new chunk representation must remain compatible with evidence resolution, citation identity, and retrieval labels. Naming a module owner is useful for future coordination, but should be agreed by the team rather than inferred from this report.

\subsection{Communication Plan}
The recommended communication process is to maintain implementation issues and benchmark findings with the code repository. Each issue should identify a concrete input, the observed behaviour, the expected contract, and the artifact needed to reproduce it. A short integration review at each work-package boundary should resolve configuration or schema changes before the next complete benchmark.

This is a proposed working process; the repository does not supply meeting minutes or evidence of a fixed weekly meeting schedule. For academic review, the supervisor can assess the documented architecture, experiment matrix, and result interpretation. Decisions affecting labels or legal scope should be recorded separately from implementation fixes so that later runs can distinguish data changes from system changes.

\section{Project Management Approach}
\subsection{Work Breakdown Structure}
The work is organized by dependency rather than by an invented calendar. Figure~\ref{fig:wbs} shows the sequence of deliverable gates. A corpus snapshot must be identifiable before evidence labels can be checked; inference artifacts must be frozen before offline judgments and statistical summaries can be audited.

\begin{figure}[H]\centering
\begin{tikzpicture}[node distance=5mm]
\node[box,text width=10cm] (a) {1. Define domain, question schema and comparison protocol};
\node[box,text width=10cm,below=of a] (b) {2. Prepare pinned source snapshot, chunks and indexes};
\node[box,text width=10cm,below=of b] (c) {3. Integrate retrieval, context, generation and verification};
\node[box,text width=10cm,below=of c] (d) {4. Validate labels and run controlled B0--B4 experiments};
\node[box,text width=10cm,below=of d] (e) {5. Score, audit, interpret uncertainty and document results};
\foreach \u/\v in {a/b,b/c,c/d,d/e}{\draw[arrow] (\u)--(\v);}
\end{tikzpicture}
\caption{Work-package dependencies and review gates. This is a delivery structure, not a record of dates.}\label{fig:wbs}
\end{figure}

\subsection{Risk Management}
\begin{table}[H]\centering\small
\caption{Material risks and controls supported by the project design.}\label{tab:risks}
\begin{tabularx}{\textwidth}{P{3.1cm}YY}\toprule
\textbf{Risk} & \textbf{Control} & \textbf{Residual limitation}\\\midrule
Missing scope or conditions & Preserve original-question filters; retain complete parent evidence where budget allows & Correct ranking may still omit required provisions\\
Fabricated citations & Source-number selection and deterministic citation construction; support verifier & Model-based verification can misjudge entailment\\
Source/version drift & Pin dataset revision and corpus identity; preserve source URLs & Snapshot does not certify current legal effect\\
Benchmark leakage & Declare development use; separate future held-out test & Existing pilot cannot be treated as unseen evaluation\\
Judge bias and truncation & Store claim judgments and errors; report valid denominators & Current quality means remain provisional\\
Slow or interrupted runs & Cache prepared corpus; atomic sample artifacts and resume checks & Model loading and local contention affect timings\\\bottomrule
\end{tabularx}
\end{table}
Risk management should prioritize observed failure modes. In this pilot, context-sufficiency false negatives and offline judge instability are more actionable than hypothetical increases in graph complexity. Their fixes require separate controls so that improving the evaluator does not silently change the underlying relevance labels.

\subsection{Quality Management}
Quality is checked at three levels. Deterministic tests enforce schemas, metadata matching, token budgets, source integrity, and failure handling. Focused regression cases exercise the interaction between retrieval scope, context packing, and citation verification. The benchmark then compares end-to-end behaviour under a shared protocol. Passing one level does not establish the next: a valid JSON response may still be unsupported, and a successful regression case is not a held-out accuracy estimate.

The latest recorded project-wide regression run contains 366 passed tests, 10 skipped tests, and 246 subtests. Two focused post-fix inference batches contain 11 B3 cases and four B0 cases, all completed, while two missing-evidence controls were rejected. These checks support the stability of the revised contracts. They are development checks, not independent experimental evidence of general legal-answering accuracy.

\subsection{Budget and Funding}
No expenditure ledger, paid-service budget, or measured energy cost is included in the project evidence. The prototype uses open-source software and locally served pretrained models. This removes the need to report an assumed commercial API bill, but does not make computation cost-free. Hardware ownership, electricity, model storage, and researcher time remain unmeasured. Future cost reporting should record these quantities together with the number of completed and judged questions.

\subsection{Change Management Process}
A meaningful change should include an explicit reason and a narrow verification case. Retrieval or orchestration changes should be checked against the original question, retained evidence, and final citation behaviour. Dataset edits require a new version and a documented explanation. A complete comparison should use a new output directory so that artifacts from different code or label versions are not mixed.

The run manifest stores configuration, code identity, model identities, question contents, and the latency protocol. The scorer stores its own identity and fingerprints. These records support an auditable chain from the selected source snapshot to a reported number. They do not replace a clean code snapshot: the report is based on the current working project and frozen artifacts, and its regeneration instructions preserve the distinction between rebuilding a document and rerunning inference.

\subsection{Closure and Evaluation}
The present milestone closes with a working benchmark cycle and a documented development result. Closure requires that all 225 expected sample and score artifacts are accounted for, report tables agree with those artifacts, and unresolved issues are stated. The next research milestone should require an independently reviewed held-out set, more complete hardware records, and calibration of the offline judge. A statistically significant advantage is not a prerequisite for an academically useful result; transparent evidence that a proposed addition did not demonstrate a benefit is itself informative.
